\chapter{Studies}\label{ch:studies}

This chapter describes the eleven studies: what each one computes and how, its options,
what its result tab shows, how it checks itself, and what it draws on the diagram. The
option tables are generated from the studies' definitions.

\begin{reftable}{@{}>{\raggedright\arraybackslash}p{4.6cm}>{\raggedright\arraybackslash}p{2.6cm}X@{}}
\toprule
\sffamily\bfseries Study & \sffamily\bfseries Menu group & \sffamily\bfseries In short \\ \midrule
Per-unit model (\code{model}) & Model & Bases and every component in per unit. \\
Admittance matrix (\code{ybus}) & Model & $Y$ of the positive, negative or zero sequence. \\
Sequence model (\code{seqmodel}) & Model & What path each element gives in each sequence. \\
S1a) Two-port matrix (\code{s1a}) & Assignment S1 & ABCD matrix between two buses. \\
S1b) Equivalent impedance (\code{s1b}) & Assignment S1 & $Z = U/I$ into the passive network. \\
S1c) Bus voltages (\code{s1c}) & Assignment S1 & Linear solve with the S1 phasor chain. \\
Bus voltages (\code{linear}) & Network & Linear solve of any network. \\
Thevenin equivalent (\code{thevenin}) & Network & $U_\text{th}$ and $Z_{\text{th},1/2/0}$ at a bus. \\
Short circuit (\code{short\_circuit}) & Network & 3ph, SLG, LL or LLG fault. \\
Load flow (\code{loadflow}) & Network & Newton--Raphson load flow. \\
Unbalanced load (\code{unbalanced\_load}) & Network & A load with different power per phase. \\
\bottomrule
\end{reftable}

\paragraph{Which part of the network a study sees.} Out-of-service elements are always
ignored. The network studies (admittance matrix, linear voltages, Thevenin, short
circuit, load flow, unbalanced load, sequence model) also keep only the \emph{energized}
part of the network: the islands that contain an external grid or a generator. A case
that disconnects part of the network simply sees the rest, and the disconnected buses are
reported as information. The S1 studies analyse the whole network as drawn.

\paragraph{Loads in the linear studies.} Every study except the load flow treats every
load as a constant impedance at its nominal voltage, which keeps the problem linear and
exact in one solve. A load set to \code{constant\_power} gets a warning in those studies.

% ===========================================================================
\section{Per-unit model}\label{sec:study-model}

The starting point of every hand calculation: the bases and each component in per unit,
with the factors that took it there.

% Generated by tools/make_manual.py. Do not edit.
\noindent This study has no options.


\paragraph{Result.}
\begin{itemize}
  \item \gui{Base values}: $U_\text{base}$, $Z_\text{base}$ and $I_\text{base}$ of every
    bus, with $S_\text{base}$ in the note;
  \item \gui{Voltage zones}: the buses that share each base voltage;
  \item \gui{Component values in per unit}: for each component, its $Z$, its shunt
    admittance (the PI half-shunt for lines, the total admittance for loads and shunts),
    its power in MVA, its nominal voltage, the factor $(U_\text{base}/U_n)^2$ and a note,
    such as \emph{2 parallel line(s), 0.297 km} or \emph{capacitive}.
\end{itemize}

% ===========================================================================
\section{Admittance matrix (Y-bus)}\label{sec:study-ybus}

The bus admittance matrix of one sequence network of the energized network, the way the
course writes it: S2a (positive sequence, source as an internal node, $9\times9$) and S3a
(zero sequence, source grounded, $8\times8$).

% Generated by tools/make_manual.py from the options of study 'ybus'. Do not edit.
\begin{optiontable}
\gui{Sequence}\newline\code{sequence} & choice & One of \code{positive}, \code{negative}, \code{zero}. \\
\gui{Sources}\newline\code{sources} & choice & internal\_node: each source behind its impedance becomes an extra node (the S2a form). grounded: its impedance goes to ground at the bus (the S3a form). exclude: leave sources out. One of \code{internal\_\allowbreak{}node}, \code{grounded}, \code{exclude}. \\
\end{optiontable}


\paragraph{Sources.} A source with zero impedance never enters the matrix, because as an
admittance it would be a short circuit. A source behind an impedance enters as follows:
\begin{description}
  \item[\code{internal\_node}] as an extra node \emph{(internal)} behind a series branch
    $Z_\text{th}$, placed after the buses;
  \item[\code{grounded}] as a shunt $1/Z_\text{th}$ to earth at its bus (the Norton form);
  \item[\code{exclude}] not at all: the passive network.
\end{description}

\paragraph{Result.} The matrix, with MATLAB's common factor in the title when there is
one (\emph{1.0e+02 *}). Below it are the \gui{Nodes} table, which lists each row with
its node, id and $U_\text{base}$, and a \gui{Checks} table with the symmetry residual
$\max|Y - Y^T|$ and the source treatment. Right-click the matrix for \gui{Copy as MATLAB}
(\cref{sec:copy-matlab}).

% ===========================================================================
\section{Sequence model}

For each element, the path it gives in the positive and the zero sequence, with its
impedance or admittance in per unit and a note. The note says, for example, that a
\code{Dyn} transformer is a shunt on its star side, that a delta load is open, or that a
grid's $Z_0$ comes from $3U/I_{k1} - 2Z_1$. The negative sequence equals the positive one
without source voltages, so it is not printed separately. When a line has no
zero-sequence data, the zero-sequence table is replaced by a warning that names it.

% Generated by tools/make_manual.py. Do not edit.
\noindent This study has no options.


% ===========================================================================
\section{S1a) Two-port matrix}\label{sec:study-s1a}

The transmission (ABCD) matrix between two buses, with the course's convention
\[
  \begin{bmatrix} U_1 \\ I_1 \end{bmatrix} =
  F \begin{bmatrix} U_2 \\ I_2 \end{bmatrix}, \qquad
  F = \begin{bmatrix} A & B \\ C & D \end{bmatrix},
\]
port~1 on the source side and port~2 on the load side. Swapping the ports gives a
different matrix.

% Generated by tools/make_manual.py from the options of study 's1a'. Do not edit.
\begin{optiontable}
\gui{From bus (port 1, source side)}\newline\code{bus\_\allowbreak{}i} & bus & Any bus of the network. \\
\gui{To bus (port 2, load side)}\newline\code{bus\_\allowbreak{}j} & bus & Any bus of the network. \\
\gui{Shunts at the ports}\newline\code{include\_\allowbreak{}at\_\allowbreak{}ports} & elements & A load or source sitting on a port bus is the termination, not part of the two-port, so it is excluded. Tick one to include it anyway. Line charging is always kept. \\
\end{optiontable}


\paragraph{The termination rule.} The two-port is the network \emph{between} the ports.
A load or source sitting on a port bus is the termination, not part of the two-port, so
it is excluded. The same holds for whatever hangs off a port without leading to the other
one, such as the rest of a feeder beyond port~2. The dialog lists the shunts on the port
buses, and you can tick one to include it anyway. An ideal source on a port can never be
included, because it would short-circuit the port. Line charging is always kept, because
it belongs to the line.

\paragraph{Two methods.} The matrix is computed twice. First by cascading the element
matrices along the path,
\[
  F_\text{series}(z) = \begin{bmatrix} 1 & z \\ 0 & 1 \end{bmatrix}, \quad
  F_\text{shunt}(y) = \begin{bmatrix} 1 & 0 \\ y & 1 \end{bmatrix}, \quad
  F_\pi = F_\text{shunt}(y_\text{sh})\,F_\text{series}(z)\,F_\text{shunt}(y_\text{sh}),
\]
and then by Kron reduction of the Y-bus onto the two ports. The relative residual between
the two is reported. It is a real check: the line charging lives inside $F_\pi$ in one
method and on the diagonal of $Y$ in the other, so counting it twice would show up as a
residual of about $10^{-5}$. In a meshed network the path is not unique and the cascade
does not exist; the result then comes from the Y-bus alone, with a warning.

\paragraph{Result.} The convention, the elements excluded at the ports and why, each
element's matrix in cascade order, $F$ itself, and the checks: $\det F$ (equal to 1 for a
reciprocal network), $|\det F - 1|$, the cascade-versus-Y-bus residual and the cascade
path.

% ===========================================================================
\section{S1b) Equivalent impedance}

The impedance $Z = U/I$ seen looking into the \emph{passive} network from one bus, with
the sources switched off. In S1, this is the factory seen from its feeding point.

% Generated by tools/make_manual.py from the options of study 's1b'. Do not edit.
\begin{optiontable}
\gui{Seen from bus}\newline\code{bus} & bus & Any bus of the network. \\
\gui{Cross-check against port}\newline\code{far\_\allowbreak{}bus} & bus & Optional: recompute through the two-port relation and report the residual. \\
\end{optiontable}


The impedance comes from Kron reduction of the Y-bus onto the bus. With a
\gui{Cross-check against port}, it is computed again through the two-port relation
$Z = (A Z_\text{term} + B)/(C Z_\text{term} + D)$, with $Z_\text{term}$ the termination at
the far port, and both residuals are reported.

\paragraph{Result.} $Z$ in per unit and in ohm, $|Z|$, $\arg Z$ and the base impedance,
then the cross-check table.

% ===========================================================================
\section{S1c) Bus voltages and the linear network}\label{sec:study-linear}

Two studies share one solve: \gui{S1c) Bus voltages} and \gui{Bus voltages (linear
network)}. They solve the network driven by its sources, with every load a constant
impedance, so one linear solve is exact:
\[
  Y U = I, \qquad
  \text{ideal sources fixing their bus voltage,}\quad
  \text{others a Norton pair } \left(\tfrac{1}{Z_\text{th}},\ \tfrac{E}{Z_\text{th}}\right).
\]
This is the pre-fault state the Thevenin and fault studies start from.

% Generated by tools/make_manual.py from the options of study 's1c'. Do not edit.
\begin{optiontable}
\gui{Report voltage at bus}\newline\code{bus} & bus & Any bus of the network. \\
\end{optiontable}


\gui{Bus voltages (linear network)} has the same option:

% Generated by tools/make_manual.py from the options of study 'linear'. Do not edit.
\begin{optiontable}
\gui{Report voltage at bus}\newline\code{bus} & bus & Any bus of the network. \\
\end{optiontable}


\paragraph{Result.}
\begin{itemize}
  \item the headline: the voltage at the chosen bus in pu, angle and kV, and the voltage
    drop from the source in percent;
  \item \gui{Phasor chain} (single-source networks): the S1 chain $U_a \to I_a \to [U_e;
    I_e]$, with $I_a = U_a/Z_I$ and $[U_e; I_e]$ from the two-port matrix;
  \item \gui{Bus voltages}: every bus in pu, angle, kV, and the current injected;
  \item \gui{Branch flows}: each branch's current, power and losses at its sending end,
    with the total losses in the note;
  \item \gui{Sources}: each source's current in pu and kA, and its power in MVA;
  \item \gui{Cross-check}: the far-bus voltage from the nodal solve against the two-port
    chain, and the source current against $U/Z$.
\end{itemize}
With more than one source, the phasor chain and the two-port cross-check are skipped,
with a warning, because $I = U/Z$ then no longer holds. The nodal solve is exact either
way.

\paragraph{Annotations.} Every bus shows its voltage in kV and in pu with its angle, and
every branch its current and sending-end power.

% ===========================================================================
\section{Thevenin equivalent}\label{sec:study-thevenin}

The equivalent seen from one bus: the open-circuit voltage and the impedance of each
sequence network.

% Generated by tools/make_manual.py from the options of study 'thevenin'. Do not edit.
\begin{optiontable}
\gui{Seen from bus}\newline\code{bus} & bus & Any bus of the network. \\
\end{optiontable}


\paragraph{How.} $U_\text{th}$ is the bus voltage of the pre-fault linear solve, with no
voltage factor. $Z_{\text{th},1}$, $Z_{\text{th},2}$ and $Z_{\text{th},0}$ are the diagonal
entries $Z_{kk}$ of the impedance matrices of the three sequence networks, with the
sources' EMFs switched off and their impedances kept. They are computed as columns of
$Z = Y^{-1}$ from one LU factorisation per sequence. A bus with no path to earth in a sequence,
such as the zero sequence behind a delta winding, has an \emph{open} (infinite) impedance
there. When a line has no zero-sequence data, $Z_{\text{th},0}$ is left out with a
warning.

\paragraph{Cross-checks.} $Z_{\text{th},1}$ is compared with the Kron reduction of the
Y-bus onto the bus. With a single source behind an impedance, $U_\text{th}$ is also
compared with the route of the MATLAB solution, $U_\text{th} = Z(k,s)\,E/Z(s,s)$. The
residuals are in the \gui{Cross-checks} table.

\paragraph{Result.} $U_\text{th}$ in pu and kV; each $Z_\text{th}$ in pu and ohm; the
three-phase fault current $I_{sc} = U_\text{th}/Z_{\text{th},1}$ in pu and kA; the
short-circuit power as the course defines it, $S_{sc} = U_\text{th}\,\conj{I_{sc}}\,
S_\text{base}$, its magnitude, and $\sqrt3\,U_\text{base}|I_{sc}|$ for comparison; and,
when $Z_0$ is known, the line-to-ground fault current $3U_\text{th}/(Z_1 + Z_2 + Z_0)$.

\paragraph{Annotations.} The bus shows $U_\text{th}$, $Z_{\text{th},1}$ and
$Z_{\text{th},0}$.

\begin{note}
A bus that holds an ideal source has $Z_\text{th} = 0$ by construction. \app\ refuses it
and asks for a finite source impedance. The default bus avoids such a bus.
\end{note}

% ===========================================================================
\section{Short circuit}\label{sec:study-short-circuit}

A shunt fault at one bus, with every bus voltage and branch current during the fault,
phase by phase.

% Generated by tools/make_manual.py from the options of study 'short_circuit'. Do not edit.
\begin{optiontable}
\gui{Faulted bus}\newline\code{bus} & bus & Any bus of the network. \\
\gui{Fault type}\newline\code{kind} & choice & 3ph: three-phase. slg: phase a to ground. ll: phases b and c. llg: phases b and c to ground. One of \code{3ph}, \code{slg}, \code{ll}, \code{llg}. \\
\gui{Fault impedance Z\_f [ohm]}\newline\code{zf\_\allowbreak{}ohm} & complex & Per faulted phase; 0 is a solid fault. \\
\gui{Ground impedance Z\_g [ohm]}\newline\code{zg\_\allowbreak{}ohm} & complex & Common path to earth of a double line-to-ground fault. Shown when \gui{Fault type} is \code{llg}. \\
\end{optiontable}


\paragraph{Fault types.} With phase $a$ as the reference phase:
\begin{description}
  \item[\code{3ph}] all three phases to a common point through $Z_f$ each;
  \item[\code{slg}] phase $a$ to earth through $Z_f$;
  \item[\code{ll}] phases $b$ and $c$ together through $Z_f$;
  \item[\code{llg}] phases $b$ and $c$ to earth, each through $Z_f$ and together through
    $Z_g$.
\end{description}

\paragraph{How.} Each fault is written as three phase equations $A\,U_{abc} + B\,I_{abc} =
0$ and solved in symmetrical components against the sequence Thevenin equivalents of the
bus (\cref{sec:shunt-solver}). For a solid single line-to-ground fault this reduces to
the textbook $I_1 = I_2 = I_0 = U_\text{th}/(Z_1 + Z_2 + Z_0 + 3Z_f)$. Only the sequences
a fault involves are built: a three-phase fault needs neither the negative nor the zero
sequence. The voltages at every other bus follow from one impedance-matrix column per
sequence, $U_n(i) = U_{\text{pre},n}(i) - Z_n(i,k)\,I_n$, and the branch currents from
those voltages. There is no voltage factor: the pre-fault state is the linear network.

\paragraph{Result.}
\begin{itemize}
  \item \gui{Fault}: bus, type, $Z_f$ (and $Z_g$) in ohm and pu, $U_\text{pre}$, the
    sequence impedances, $I_1$, $I_2$, $I_0$ and $I_f$ in pu and kA, the earth current
    $3I_0$ for faults to earth, the fault power $S_f = U_\text{th}\conj{I_f}S_\text{base}$
    for a three-phase fault, and $\sqrt3\,U_\text{base}|I_f|$;
  \item \gui{At the faulted bus}: the phase currents and line-to-neutral voltages;
  \item \gui{Bus voltages during the fault}: $U_a$, $U_b$, $U_c$ of every bus in pu and kV;
  \item \gui{Branch currents during the fault}: each branch's phase currents and its
    largest current in kA;
  \item the Thevenin cross-checks.
\end{itemize}

\paragraph{Annotations.} The faulted bus shows $I_f$ in kA and its three phase voltages.
Other buses show their voltage, or their three phase voltages when the fault unbalances
them, and every branch shows its largest current.

\begin{tip}
For a quick three-phase fault, right-click a bus and choose \gui{Three-phase short circuit
here}.
\end{tip}

% ===========================================================================
\section{Load flow (Newton--Raphson)}\label{sec:study-loadflow}

The steady state of the network with its generators, voltage controllers and
constant-power loads: the problem of Assignment LF.

% Generated by tools/make_manual.py from the options of study 'loadflow'. Do not edit.
\begin{optiontable}
\gui{Load model}\newline\code{load\_\allowbreak{}model} & choice & as\_set: each load's own model. constant\_power: every load is a PQ injection. constant\_impedance: every load is an admittance at its nominal voltage. One of \code{as\_\allowbreak{}set}, \code{constant\_\allowbreak{}power}, \code{constant\_\allowbreak{}impedance}. \\
\gui{Shunt model}\newline\code{shunt\_\allowbreak{}model} & choice & as\_set: each shunt's own model. constant\_q: the shunt injects its rated Q whatever the voltage. constant\_b: the rating fixes B at U nominal, so it delivers B U\textasciicircum{}2. One of \code{as\_\allowbreak{}set}, \code{constant\_\allowbreak{}q}, \code{constant\_\allowbreak{}b}. \\
\gui{Enforce reactive limits}\newline\code{q\_\allowbreak{}limits} & check box & A PU bus whose generators or SVC would leave their Q limits becomes a PQ bus held at the limit. \\
\gui{Tolerance [pu]}\newline\code{tolerance} & number & Largest power mismatch accepted. At least 0. \\
\gui{Maximum iterations}\newline\code{max\_\allowbreak{}iter} & integer & At least 1. \\
\end{optiontable}


\paragraph{Bus types.} \app\ derives each bus's type from what is connected, and names
it as the course does:
\begin{description}
  \item[Slack] a bus with a slack generator or an external grid with \gui{Slack} checked.
    Its voltage and angle are given. Exactly one per island is required.
  \item[PU] a bus with a PV generator or an SVC. Its active power and voltage magnitude
    are given.
  \item[PQ] every other bus. Its active and reactive power are given (zero at a bus with
    nothing connected).
\end{description}

\paragraph{What enters where.} Constant-impedance loads and constant-B shunts are part of
the admittance matrix. Constant-power loads, constant-Q shunts, PQ generators and the
fixed $P$ of PV generators are \emph{injections}, with generation counted positive.
External grids enter only as slack buses; their impedance is short-circuit data.

\paragraph{The method.} Polar Newton--Raphson from a flat start (1~pu and $0^\circ$
except at the given voltages). The mismatches are
\[
  \Delta P_k = P_k^\text{spec} - P_k(\theta, U) \ \ (\text{PU and PQ buses}),\qquad
  \Delta Q_k = Q_k^\text{spec} - Q_k(\theta, U) \ \ (\text{PQ buses}),
\]
and each iteration solves for the angle and relative magnitude corrections
\[
  \begin{bmatrix} \Delta P \\ \Delta Q \end{bmatrix} =
  \begin{bmatrix} H & N \\ J & L \end{bmatrix}
  \begin{bmatrix} \Delta\theta \\ \Delta U/U \end{bmatrix}, \qquad
  \theta \mathrel{+}= \Delta\theta, \quad U \mathrel{*}= 1 + \Delta U/U,
\]
with $H = \partial P/\partial\theta$, $N = U\,\partial P/\partial U$, $J = \partial
Q/\partial\theta$ and $L = U\,\partial Q/\partial U$ (the element formulas are in
\cref{sec:nr}). The iteration stops when the largest mismatch falls below the tolerance
($10^{-10}$~pu by default), usually in four or five iterations.

\paragraph{Reactive limits.} With \gui{Enforce reactive limits} checked, a PU bus whose
generators or SVC would need more than their combined \gui{Q max} (or less than \gui{Q
min}) becomes a PQ bus, held at that limit, and the flow is solved again from the current
state. A warning names the bus. A bus that reaches a limit stays PQ: there is no
switching back in this version.

\paragraph{Result.}
\begin{itemize}
  \item the headline: convergence, iterations, the total losses and the slack's $P$ and
    $Q$;
  \item \gui{Bus types and voltages}: every bus's type, $|U|$ in pu and kV, angle, and the
    net $P$ and $Q$ it injects;
  \item \gui{Generation and voltage control}: each generator, external grid and SVC with
    its bus type, $P$ and $Q$. An SVC also reports $B = Q/U^2$ in pu and siemens, and a
    limited device says which limit holds it;
  \item \gui{Loads}: each load's model and the $P$ and $Q$ it draws at the solved voltage.
    A constant-impedance load draws $S_n\,U^2$;
  \item \gui{Shunts}: each shunt's model, the $Q$ it delivers and its $B$. A constant-B
    shunt delivers $B\,U^2$, which is why the LF variant's bank at bus~5 gives 8.3648~MVAr
    rather than its rated 8.2~MVAr;
  \item \gui{Branch flows}: $P$ and $Q$ into each branch at both ends, its active and
    reactive losses, and its largest current in kA;
  \item \gui{Losses and power balance}: $P_\text{loss}$ as the sum of the branch losses,
    $Q_\text{loss}$, the generation and the load, and the two balance residuals, which
    close at the tolerance;
  \item \gui{Newton-Raphson iterations}: for each iteration, the largest $|\Delta P|$ and
    $|\Delta Q|$, the bus where the larger one occurs, and the largest step taken;
  \item \gui{Jacobian [H N; J L] at the flat start}, with every row and column labelled.
    The Jacobian at the solution and the load-flow Y-bus are in the result too, and they
    come out with \gui{Export results} as JSON.
\end{itemize}

\paragraph{When it does not converge.} The study does not stop with an error. It shows
the last iterate with a warning, \emph{``The load flow did not converge: largest mismatch
\dots\ after 20 iterations. The values below are the last iterate, not a solution.''},
and keeps the iteration log, so you can see where it went wrong. Typical causes are a load
far beyond what the network can carry, or an island whose only voltage reference is
missing.

\paragraph{Annotations.} Every bus shows $|U|\angle\theta$, its type and its voltage in kV.
Voltages outside 0.95--1.05~pu are drawn in the warning colour. Every branch shows its
sending-end $P$ and $Q$ and its losses, every generator its $P$ and $Q$, and every SVC
its $Q$ and $B$.

% ===========================================================================
\section{Unbalanced load}\label{sec:study-unbalanced}

A load that draws a different power on each phase, solved with the sequence networks.
Assignment S3e disconnects half of phase $a$ of load LDe: factors 0.5, 1, 1.

% Generated by tools/make_manual.py from the options of study 'unbalanced_load'. Do not edit.
\begin{optiontable}
\gui{Load}\newline\code{element} & a load &  \\
\gui{Phase a factor}\newline\code{factor\_\allowbreak{}a} & number & Share of the load's phase-a power still connected; 0.5 is half of it disconnected. At least 0. \\
\gui{Phase b factor}\newline\code{factor\_\allowbreak{}b} & number & At least 0. \\
\gui{Phase c factor}\newline\code{factor\_\allowbreak{}c} & number & At least 0. \\
\gui{Report voltages at bus}\newline\code{bus} & bus & Any bus of the network. \\
\end{optiontable}


\paragraph{How.} The load is taken out of the sequence networks, which then give the
voltage at its bus without it and the sequence Thevenin impedances there. The load's
phase admittance matrix $Y_\text{ph}$, built from its balanced admittance, its connection
and the phase factors, closes the problem, solved like a fault with $A = Y_\text{ph}$ and
$B = -\mathbf{I}$ (\cref{sec:shunt-solver}). For a grounded star, $Y_\text{ph}$ is
diagonal; a neutral impedance couples the phases; an ungrounded star and a delta carry no
zero-sequence current. For a delta, the factors apply to the branches $ab$, $bc$ and
$ca$. The voltages elsewhere follow from the impedance-matrix columns.

\paragraph{Result.} The load (connection, factors, balanced impedance, voltage without it,
sequence impedances at its bus); the matrices $Z_\text{ph}$ and $Z_s = T^{-1}Z_\text{ph}T$;
the sequence and phase currents it draws and the phase voltages at its bus; the voltages
at the report bus (pre-disturbance, sequence and phase); and every bus's three phase
voltages.

\paragraph{Annotations.} Every bus shows its three phase voltages in pu. The load's bus
and the report bus are highlighted.
